\documentclass[a4paper,dvips]{slides}
\usepackage[francais,english]{babel}
\usepackage{ifthen}
\usepackage{epsfig}
\usepackage{xspace}

\usepackage{Sty/fig-cap}
\usepackage{Sty/simple}

%%\onlyslides{1-999}

\renewcommand{\arule}[1]{\emph{\mbox{#1}}} % CA rule
\renewcommand{\asquare}[1]{\ensuremath{#1^{2}}}
\newcommand{\agen}[1]{gen.\,#1}
\newcommand{\agencnt}[1]{#1\,gen.}
\newcommand{\aplane}[1]{bit #1}

\newenvironment{mynote}{\begin{note}\begin{scriptsize}}{\end{scriptsize}\end{note}}

\newcommand{\mytitle}[1]{\begin{center}\begin{large}#1\end{large}\end{center}}
\newcommand{\mystitle}[1]{\begin{center}#1\end{center}}

\newcommand{\myslide}[3]{%
  \begin{slide}
    \mystitle{#1}
    \vspace{\bigskipamount}
    #2
  \end{slide}
  \begin{mynote}
    #3
  \end{mynote}
}

\sloppy
\selectlanguage{francais}

\begin{document}

\addtolength{\textheight}{\topmargin}
\addtolength{\textheight}{\headheight}
\addtolength{\textheight}{\headsep}
\addtolength{\textheight}{\footskip}

\setlength{\topmargin}{0pt}
\setlength{\headheight}{0pt}
\setlength{\headsep}{0pt}
\setlength{\footskip}{0pt}

%%
%%% INTRO
%%

%%
%% TALK TITLE
%%
\begin{center}
  \begin{Large}
    Exploration et visualisation\\
    d'espaces du calcul cellulaire
  \end{Large}
  \bigskip
  \begin{large}
    Thierry Delamare\\
  \end{large}
  \bigskip
  \bigskip
  DEA d'intelligence artificielle\\
  Universit\'e Paris-8 et Paris-13\\
  \bigskip
  \today
  \bigskip
\end{center}
\bigskip
\bigskip
\bigskip
\mySfigure{twothirds}{dump}{qmean-9-384}%
{\arule{Qmean} \asquare{512} \agen{384} \aplane{3}}

%%
%% TALK STRUCTURE
%%
%%\begin{slide}
%%  \mytitle{Structure de l'expos\'e\,:}
%%  
%%  Les r\`egles utilis\'ees.
%%
%%  Les architectures des moteurs r\'ealis\'es.
%%
%%  Les visualisations r\'ealis\'ees.
%%
%%  D\`es r\`egles aux espaces de r\`egles.
%%  
%%  Du cellulaire au SIMD.
%%\end{slide}

\begin{slide}
  \begin{center}
    \begin{Large}
      Th\`emes de l'expos\'e\,:
    \end{Large}
  \end{center}
  
  R\`egles et espaces de r\`egles.
  
  Architecture des moteurs de calcul et \'equi\-valence cellulaires-SIMD.
  
  Exploration et visualisation de trajectoires.
\end{slide}

%%
%%% RULES USED
%%
\begin{slide}
  \mytitle{Trois r\`egles additives\\
    et leurs tables secondaires\,:}
  \begin{tabbing}
    XXXXXX\=\kill
    \arule{Life}\>\texttt{00U100000}\\
    \arule{Anneal}\>\texttt{0000101111}\\
    \arule{Qmean}\>\texttt{mask = 3; match = 0; delta = -1;}\\
  \end{tabbing}
  \begin{scriptsize}
    \vspace{-\bigskipamount}
    \begin{tabbing}
      XX\=\kill
      \>\texttt{/* Qmean secondary table computation */}\\
      \>\texttt{/* Construct Mean */}\\
      \>\texttt{m = (SUM + C) / 9;}\\
      \>\texttt{/* Destruct it */}\\
      \>\texttt{return~m~\&~mask~==~match~?~m~:~m~+~delta;}
    \end{tabbing}
  \end{scriptsize}
  \bigskip
  \mySthreefigure{full}{dump}{hor}{life-00}{anneal-00}{qmean-9-300}%
  {\begin{scriptsize}
    \arule{Life} \asquare{128} \agen{128},
    \arule{Anneal} \asquare{256} \agen{256},\\
    \arule{Qmean} \asquare{512} \agen{300} \aplane{3}
    \end{scriptsize}}
\end{slide}

%%\begin{mynote}
%%  THE RULES USED
%%
%%  Toutes ces r\`egles travaillent sur des espaces de dimension 2, en
%%  voisinnage de Moore.
%%
%%  Toutes sont des r\`egles additives.\\
%%  Leurs tables de transitions peuvent \`etre exprim\'e par une table
%%  secondaire, dont les entr\'ees sont fonction de la somme du voisinnage.
%%
%%  \arule{Life} et \arule{Anneal} sont des r\`egle binaire (le nombre
%%  d'\'etats par cellule est 2).
%%
%%  Le nombre d'\'etats de \arule{Qmean} est 256.
%%
%%  La r\`egle \arule{Anneal} est un vote majoritaire perturb\'e.
%%  Vichniac la n\^ome \arule{M46789}, le nom \arule{Anneal} est du \`a
%%  Margolus et Toffoli.
%%
%%  La r\`egle \arule{Qmean} \`a une taille virtuelle de $2^{72}$ bytes.
%%  Sa table secondaire \`a une taille de $2^{12}$ bytes.
%%  
%%  Le moteur cellulaire utilis\'e pour \arule{Qmean} fournit un
%%  voisinnage virtuel compos\'e  de la cellule centrale (C) et de la
%%  somme des huit voisins (SUM).
%%
%%  J'ais d\'efini \arule{Qmean} pour tester ce moteur cellulaire.
%%
%%  LIFE BITMAP
%%
%%  Chaque visualisation identifie la r\`egle utilis\'e, la taille de
%%  l'espace et la g\'en\'eration ou la s\'equence de g\'en\'eration visualis\'ee.
%%  
%%  Cette visualisation montre une configuration cellulaire
%%  situ\'e sur une trajectoire.
%%
%%  La configuration cellulaire initiale des trajectoires visualis\'e est
%%  par d\'efaut une chaine de bit obtenu par tirage \`a pile ou face.
%%
%%  ANNEAL BITMAP
%%
%%  Les configurations cellulaire des chaines de bits dot\'e d'une
%%  structure d'interconnection locale.
%%
%%  Le choix d'une structure d'interconnection fixe un des param\`etre
%%  (une des dimensions) du calcul cellulaire.
%%
%%  Un autre param\`etre est donn\'e par la r\`egle ou fonction de transition.
%%  A structure d'interconnexion constante le nombre de r\`egle est fini.
%%
%%  Le dernier param\`etre est fix\'e par le choix d'une configuration
%%  cellulaire initiale.
%%\end{mynote}

%%
%%% CELLULAR MOTORS
%%

\begin{slide}
  \mytitle{Les moteurs\,:}
  \myStwofigureA{full}{fig}{hor}{apply-1}{func-4}%
  {Pipeline d'application de \arule{Life} et \arule{Anneal}\\
    et deux tables non disjointes}
  \mySfigure{full}{fig}{func-7}%
  {Fonction d'application du voisinage de \arule{Qmean}}
\end{slide}

%%\begin{mynote}
%%  MOTOR PIPELINE
%%
%%  Cette figure montre comment pipeliner la fonction de mapping pour un
%%  voisinage de Moore, r\'eduisant l'obtention d'un nouvel index \`a
%%  1~shift et 3~copies, \`a la place des 9~copies n\'ecessaires en
%%  l'absence de pipeline.
%%
%%  TWO TABLE
%%
%%  Mise en {\oe}uvre de deux AC coordonn\'es par l'interm\'ediaire de deux
%%  r\'eseaux d'interconnexions non disjoints partageant un m\^eme espace
%%  d'exploitation, permettant la transmission d'information d'un plan
%%  cellulaire \`a l'autre.
%%
%%  SPECIAL NEIGHBOURHOOD
%%
%%  Une fonction d'application pour un voisinage de Moore o\`u l'addition
%%  des cellules sources remplace leur concat\'enation. %
%%  Les connexions barr\'ees sont activ\'ees au cycle~2, et les connexions
%%  doublement barr\'ees au cycle~3 du pipeline.
%%\end{mynote}

%%
%% ANNEAL SPECIAL SECTION
%%

\begin{slide}
  \mystitle{Leur utilisation pour l'exploration}
  \mySfigure{full}{plot}{anneal-04}%
  {Populations et tailles des fronti\`eres normalis\'ees pour 4~espaces
    initiaux de taille crois\-sante}%
  \mySfigure{full}{plot}{anneal-02}%
  {\'Evolution de la population de~\arule{Anneal} lors
    d'un changement du nombre d'\^iles}%
\end{slide}

%%\begin{mynote}
%%  La taille des fronti\`eres est une approximation par la diff\'erence bit
%%  \`a bit entre l'espace de g\'en\'eration~$n$ et~$n-1$.
%%
%%  Cette figure permet de comprendre l'existence des points
%%  d'inflexions de la courbe d'\'evolution de anneal.

%%  R\'eductions int\'egrales normalis\'ees des populations des trajectoires
%%  associ\'ees \`a 4~espaces initiaux de tailles respectives:~$2^{128}$,
%%  $2^{256}$, $2^{512}$ et $2^{1024}$. %
%%  Le plot de gauche pr\'esente le logarithme du plot de droite.
%%\end{mynote}

%%
%%% VISUALISATIONS DONE
%%

%%
%% LIFE XOR SEQUENCE AND SUM CONTOUR PLOT
%%
\begin{slide}
  \mytitle{Visualisations\,:}
  \mystitle{Diff\'erences et accumulations}
  \mySfigure{full}{dump}{life-time-xor}%
  {\arule{Life} \asquare{64} \agencnt{6}}
  \myStwofigure{full}{plot}{hor}{lifex-sum-1}{lifex-sum-2}%
  {\arule{Life} \asquare{64} \agencnt{512}}
\end{slide}

%%\begin{mynote}
%%  LIFE XOR SEQUENCE
%%
%%  qualification des visualisation par le type de r\'eduction op\'er\'e.
%%
%%  Modes de r\'eduction:\\
%%  Extraction\\
%%  Accumulation\\
%%  Diff\'erence\\
%%  Relativiste.
%%  
%%  S\'equence de r\'eductions diff\'erentielles.
%%
%%  LIFE AND LIFE XOR SUM CONTOUR PLOT
%%
%%  Visualisation par countours\\
%%  une readuction par accumulation\\
%%  une r\'eduction par accumulation de diff\'erence.
%%
%%  Filtre de mise en valeur des propri\'et\'es de la trajectoire\\
%%  Par exemple, possible r\'ev\'elateur d'invariant par changement
%%  d'echelle d'espace pour life.
%%
%%  R\'ev\'elateur de l'activit\'e au fronti\`eres pour anneal.
%%\end{mynote}

%%
%% LIFE AND LIFE XOR, 3D VIS
%%
\begin{slide}
  \mystitle{Espace-temps}
  \mySfigure{full}{dump}{life-3d-New}%
  {\arule{Life} \asquare{64} \agencnt{512}}
\end{slide}

%%\begin{mynote}
%%  extraxion
%%
%%  diff\'erence + extraxion
%%\end{mynote}

  \mySthreefigure{full}{dump}{hor}{tsum-128x3-diag-1}{anneal-256-rel}{life-256x3-relmov}%
  {foobar}
  \mySthreefigure{full}{dump}{hor}{tsum-128-faces}{tsum-128-trans-byte}{tsum-128-trans-bits}%
  {barfoo}
  \myStwofigureB{full}{dump}{hor}{tsum-ray-one}{anneal-sum}{foo}{plot}
\begin{slide}
  \mySfigure{full}{dump}{slide-2}{foo}
\end{slide}

%%
%% LIFE FIXED AND MOBILE RELATIVISTIC PROJECTION
%% ANNEAL FIXED RELATIVISTIC PROJECTION
%%
\begin{slide}
  \mystitle{Projections relativistes \`a cellule de r\'ef\'erence fixe et mobile}
  \myStwofigureA{full}{dump}{hor}{life-256x3-rel}{life-256x3-relmov}%
  {\arule{life} \acube{256}}
  %%
  \mystitle{Projection relativiste sur les 3~axes d'un cube d'espace-temps}
  \mySfigure{full}{dump}{anneal-256x3-rel}%
  {\arule{Anneal} \acube{256}}  
\end{slide}

%%\begin{mynote}
%%  LIFE FIXED RELATIVISTIC PROJECTION
%%
%%  Extraction relativiste.\\
%%  R\'ev\'elateur dense de la structure globale de l'espace-temps.
%%
%%  LIFE MOBILE RELATIVISTIC PROJECTION
%%
%%  ANNEAL FIXED RELATIVISTIC PROJECTION
%%\end{mynote}

%%
%%% FROM RULES TO RULE SPACE
%%
\begin{slide}
  \mytitle{Des r\`egles aux espaces de r\`egles\,:}
  \begin{footnotesize}
    Soit une r\`egle $R$ d'un AC de $k = 2^n$ \'etats,\\
    $d$ dimensions et de rayon de voisinage $r$.\\
    $N = k^{dr+1} = 2^b$ est le nombres d'entr\'ees de la table.
  \end{footnotesize}
  \begin{eqnarray}
    \footnotesize
    \lambda & = & (N-\#0)/N \nonumber \\
    \small
    {\rm flip} & = & \frac{1}{N}\sum_{i=0}^{N-1}{\rm
      diff}(R_{i}-R_{(i+1){\rm mod}N}) \nonumber \\
    {\rm diff}(x) & = & \left\lbrace
    \begin{array}{l}
      1\ {\rm if}\ x \ne 0\\ 0\ {\rm if}\ x = 0
    \end{array} \right. \nonumber
  \end{eqnarray}
  \begin{footnotesize}
    Pour un AC de $k = 2$ \'etats les $b = dr+1$ bits du voisinage donnent
    un vecteur $E$ de $b$ valeurs.
  \end{footnotesize}
  \begin{eqnarray}
    E_{b}& = &\frac{1}{N} \sum_{i = 0}^{N/2 - 1} R_{e_{ib}} -
    R_{e_{ib} + 2^b} \nonumber \\
    e_{ib}& = &2i + i \bmod 2^b \nonumber \\
    \hat{E}& = &\frac{1}{b} \sum_{i = 0}^{b} E_b \nonumber
  \end{eqnarray}
\end{slide}

%%
%% ENCHASSEMENT
%%
\begin{slide}
  \mystitle{Valeurs de \alambda et de l'ench\^assement pour les
    $\frac{0\cdots 500}{59049}$ AC additifs \texttt{01U}}
  \mySfigure{full}{plot}{E-lambda-1}%
  {Valeurs de \protect\alambda et $E$}
  \begin{eqnarray}
    E_{life} & = & {0.617\cdots 0.890\cdots 0.617} \nonumber \\
    \hat{E}_{life} & = & 0.647569 \nonumber \\
    \lambda_{life} & = & 0.273438 \nonumber
  \end{eqnarray}
\end{slide}

%%\begin{mynote}
%%  Nous avons calcul\'e les valeurs de \alambda et \En{} pour un
%%  sous-espace \'echantillon de l'espace des r\`egles binaires \`a
%%  voisinage de Moore.
%%
%%  Le sous espace est constitu\'e des r\`egles additives, centre inclu,
%%  dont la table secondaire supporte les fonctions constantes
%%  \aquote{z\'ero} et \aquote{un} et la fonction \aquote{unchanged}.
%%  
%%  L'espace contient donc $3^{10} = 59049$ r\`egles.
%%
%%  Les valeurs plot\'ees pour \En sont:~$En_{centre}$, $En_{rayon}$
%%  et la valeur moyenne de \En sur tous les  voisins.
%%
%%  Le pic montre la position de \arule{Anneal} dans le sous-espace de
%%  r\`egles.
%%\end{mynote}

%%
%% FROM CA TO SIMD
%%
\begin{slide}
  \mytitle{Du cellulaire au SIMD\,:}
  Transformation d'une architecture SIMD en AC \`a voisinage sp\'ecial pour la
  m\'emoire et les instruction.
  \myStwofigureA{full}{fig}{hor}{gap-pe}{gap}%
  {Structure d'un PE GAPP\\
    et fonction d'application de la table \arule{Gapp}}
  Transformation de la table de transition d'un AC en DFA compilable en
  algoritme SIMD.
  \myStwofigureA{full}{fig}{hor}{life-tree}{life-dfa}%
  {Une partie de l'arbre de d\'ecodage de \arule{Life}\\
    et Le DFA de \arule{Life}}
\end{slide}

%%\begin{mynote}
%%  PE GAPP
%%
%%  Reproduite d'apr\`es NCR85a sans faire figurer le registre CM
%%
%%  PE GAPP TABLE
%%
%%  Les fl\`eches figurant sur le sch\'ema indique les chemins de donn\'ees des
%%  bits constituant l'index de la table \arule{gapp}. %
%%  Les conditions d'affectation et les temps de cycles internes \`a la
%%  fonction d'application sont indiqu\'es \`a droite du sch\'ema.
%%\end{mynote}

%%
%%% CONCLUSION
%%
\begin{slide}
  \begin{center}
    \begin{Large}
      Conclusion
    \end{Large}
  \end{center}
  \mytitle{Propositions pour une \'etude pratique des espaces du calcul
    cellulaire\,:}

  Un usage g\'en\'eralis\'e de table de transitions.
  \begin{itemize}
  \item Des moteurs rapides pour le calcul cellulaire sur machine
    s\'equentielle.
  \item Des transformations cellulaire-SIMD.
  \end{itemize}
  Des m\'ethodes de visualisations d'espace-temps.

  Une mesure de l'ench\^assement des tables de transitions.
\end{slide}

%%
%% TUBE-WORM MOSAIC
%%
\begin{slide}
  \mystitle{Huit variations sur les param\`etres statiques d'une r\`egle}
  \mySfigure{full}{dump}{tube-worm-mix}%
  {\arule{Tube-Worm} \asquare{256}}
\end{slide}

%%\begin{mynote}
%%  TUBE-WORM MOSAIC
%%  
%%  Une composition construite avec 8~s\'equences~(de haut en bas), de
%%  8~espaces chacune~(de gauche \`a droite) de la r\`egle appliqu\'ee sur le
%%  m\^eme espace initial.
%%  
%%  L'axe vertical est celui de la variation dans l'espace des r\`egles,
%%  via une variation des param\`etres statique de tube-worm.
%%\end{mynote}

%%
%%
%% ANNEAL 3D
%%
\begin{slide}
  \mystitle{R\'eduction implicite par extraction}
  \mySfigure{full}{dump}{anneal-3d}{\arule{Anneal} \acube{512}}
\end{slide}

%%\begin{mynote}
%%  ANNEAL 3D
%%
%%  R\'eduction par extraction implicite\\
%%  par projection solide d'espace-temps 3-D.
%%
%%  R\'ev\'elateur des textures et des structures construite de
%%  l'espace-temps.
%%\end{mynote}

%%
%% ANNEAL SUM ELEVATION SURFACE
%%
\begin{slide}
  \mystitle{Surface d'\'el\'evation d'une r\'eduction explicite\,: $\sum_{i=0}^{n-1}t_{i}xy$}
  \mySfigure{full}{plot}{anneal-sum}%
  {\arule{Anneal} \asquare{64} \agencnt{400}}
\end{slide}

%%\begin{mynote}
%%  ANNEAL SUM ELEVATION SURFACE
%%
%%  R\'eduction par accumulation explicite.
%%
%%  + R\'eduction par extraction implicite.
%%
%%  La r\'eduction se fait par projection sur un plan de sommation
%%  parall\`element \`a l'axe du temps.
%%
%%  
%%\end{mynote}

%%
%% QMEAN RAY-TRACED 
%%
\begin{slide}
  \mystitle{Visualisation par lancer de rayons\\
    sur iso-surface d'espace-temps}
  \mySfigure{full}{dump}{tsum-ray-one}%
  {\arule{Qmean} \acube{128}}
\end{slide}

%%\begin{mynote}
%%  QMEAN RAY-TRACED
%%
%%  Extraction double:\\
%%  iso-surface explicite et implicite de projection 3D -> 2D
%%
%%  espace initial = gen 8192 de \arule{Anneal}.
%%\end{mynote}

%%
%% QMEAN SIX FACES OF SPACE-TIME AND QMEAN CELL TRANSPARENCY
%%
\begin{slide}
  \mystitle{Faces de cube d'espace-temps et Transparence utilisant
    tout l'espace de valeur d'une cellule}
  \myStwofigure{full}{dump}{vert}{tsum-128x3-faces}{tsum-128x3-trans-byte}%
  {\arule{Qmean} \acube{128}}
\end{slide}

%%\begin{mynote}
%%  QMEAN SIX FACES OF SPACE-TIME
%%  
%%  Ce sont des trances d'espace-temps.
%%  
%%  QMEAN CELL TRANSPARENCY
%%  
%%  R\'eduction par accumulation.
%%\end{mynote}

%%
%% QMEAN BIT-PLANES TRANSPARENCY
%%
\begin{slide}
  \mySfigure{full}{dump}{tsum-128x3-trans-bits}%
  {\arule{Qmean} \acube{128}}
\end{slide}

%%\begin{mynote}
%%  QMEAN BIT-PLANES TRANSPARENCY
%%
%%  Extraction (tranche d'etats) + Accumulation
%%
%%  Transparence \'eclat\'ee par plan de bits.
%%
%%  Les cinq bits de poids fort de la r\`egle sont utilis\'es respectivement
%%  pour r\'ealiser une transparence sur les trois axes de l'espace-temps.
%%\end{mynote}

%%
%% QMEAN SLICE SEQUENCE
%%
\begin{slide}
  \mystitle{S\'equence de tranches d'espace-temps}
  \mySfigure{full}{dump}{tsum-128x3-diag}%
  {\arule{Qmean} \acube{128}}
\end{slide}

%%\begin{mynote}
%%  QMEAN SLICE SEQUENCE
%%  
%%  R\'eduction par Extraction.
%%  
%%  Tranche non perpendiculaire \`a l'axe du temps.
%%  
%%  R\'ealis\'es en g\'eom\'etrie discr\`ete (distance de Manhattan).
%%\end{mynote}

%%% XU
%%
\begin{slide}
  \mystitle{Un AC irr\'eductible \`a sa table}
  \mySfigure{full}{dump}{xu-photo-1}%
  {La dynamique de l'AC \arule{Xu}}%
\end{slide}

%%\begin{mynote}
%%  Un cube d'espace-temps d'ar\`ete~128, rendu par ray-tracing
%%  d'iso-sur\-face. %
%%\end{mynote}

%%
%%
%%
\begin{slide}
  \mySfigure{full}{dump}{xu-photo-1-top}%
  {Un rendu d'une tranche des valeurs du dernier
    plan de la figure pr\'ec\'edente}%
\end{slide}

%%
%%
%%
\begin{slide}
  \mySfigure{full}{dump}{xu-Cfilm-1-1}%
  {D'autres points de vues sur la dynamique de \arule{Xu}}%
\end{slide}

%%
%%
%%
\begin{slide}
  \mySfigure{full}{dump}{xu-Cfilm-1-2}%
  {Et un \'epluchage par iso-surface}%
\end{slide}

%%\begin{mynote}
%%  {R\'ealis\'e en variant lin\'eairement la valeur de seuil d'extraction d'iso-sur\-face}
%%\end{mynote}

\end{document}

%% Local Variables: 
%% mode: auto-fill
%% iso-accents-minor-mode: nil
%% TeX-master: t
%% Local IspellParsing: "latex-mode ~latin1"
%% Local IspellPersDict: "~/.ispell_lisp"
%% LocalWords: "paper dvips slides francais english babel ifthen epsfig xspace"
%% LocalWords: "Sty fig-cap gen mynote pt Thierry Delamare DEA twothirds dump XX"
%% LocalWords: "qmean SIMD XXXXXX Anneal mask secondary computation Construct it"
%% LocalWords: "Mean SUM Destruct return hor anneal apply func santes lifex-sum"
%% LocalWords: "life-time-xor lifex-sum Espace-temps d-New relativiste rel tsum"
%% LocalWords: "relmov d'espace-temps anneal-sum iso-surface tsum-ray-one tsum"
%% LocalWords: "trans-byte trans-bits diag E-lambda tube-worm-mix Tube-Worm PE"
%% LocalWords: "gap-pe GAPP gap gapp life-tree life-dfa DFA cellulaire-SIMD l'AC"
%% LocalWords: "xu-photo xu xu-Cfilm \'equi"
%% End: 
